In this small simulation study, replacing a fitted single-mode pendulum or spring-mass model by a learned slosh predictor inside the same time-minimising MPPI controller shortened spill-free transport of a brim-full container on the nominal liquid: the pendulum controller took \rhval{cmp/main/pendulum-mpc/nominal/transport_time_s/inv_ratio:3} times as long, with part of that gap attributable to the margin selection. The learned model was not the best option everywhere: a fitted two-mode spring-mass model was slightly faster on the liquid it was identified on, open-loop input shaping was fastest of all there, and at their deployed margins the slower pendulum controller spilled less on held-out liquids with pushes. What the learned model with a history encoder did provide was the better time--spill trade-off on liquids it had not been tuned for, where input shaping and the two-mode model failed. With a single linear slosh mode the learned model offered nothing over the pendulum. A liquid sensor of any of the three kinds was far more important than the choice among them.

For a system that must be the fastest spill-free transporter, these results suggest three next steps: a solver that is closer to time-optimal, since the optimiser cost us as much as the model; a margin that adapts to the inferred liquid rather than one tuned on a nominal liquid; and a test against a fluid simulation or a real container, where the open question is whether the effects a learned model can absorb are large enough to matter.
